\documentclass[11pt]{article}
\usepackage[margin=0.75in]{geometry}
\usepackage{amsmath,amssymb,amsthm}
\usepackage{enumerate}
\usepackage{booktabs}
\usepackage{graphicx}
\usepackage{microtype}
\usepackage{parskip}
\usepackage[hidelinks]{hyperref}
\allowdisplaybreaks

% Common mathematical notation.
\newcommand{\N}{\mathbb{N}}
\newcommand{\Z}{\mathbb{Z}}
\newcommand{\Q}{\mathbb{Q}}
\newcommand{\R}{\mathbb{R}}
\newcommand{\C}{\mathbb{C}}
\newcommand{\F}{\mathbb{F}}

% PeerFS-specific notation.
\newcommand{\File}{\text{File}}
\newcommand{\Path}{\text{Path}}
\newcommand{\Tree}{\text{Tree}}
\DeclareMathOperator{\fs}{fs}
\DeclareMathOperator{\dir}{dir}
\DeclareMathOperator{\sub}{sub}
\DeclareMathOperator{\Reg}{Reg}
\DeclareMathOperator{\canon}{canon}
\DeclareMathOperator{\tgt}{tgt}
\DeclareMathOperator{\dom}{dom}
\newcommand{\Name}{\text{Name}}
\newcommand{\Mode}{\text{Mode}}
\newcommand{\Rd}{\mathsf{r}}
\newcommand{\RW}{\mathsf{rw}}
\newcommand{\PF}[1]{\textsf{PF#1}}

% Theorem environments.
\newtheoremstyle{definitionbreak}
    {\topsep}
    {\topsep}
    {\normalfont}
    {0pt}
    {\bfseries}
    {.}
    {\newline}
    {}

\newtheoremstyle{plainbreak}
    {\topsep}
    {\topsep}
    {\itshape}
    {0pt}
    {\bfseries}
    {.}
    {\newline}
    {}

\newtheoremstyle{remarkbreak}
    {0.5\topsep}
    {0.5\topsep}
    {\normalfont}
    {0pt}
    {\itshape}
    {.}
    {\newline}
    {}

\theoremstyle{definitionbreak}
\newtheorem{definition}{Definition}[section]

\theoremstyle{plainbreak}
\newtheorem{theorem}[definition]{Theorem}
\newtheorem{proposition}[definition]{Proposition}
\newtheorem{lemma}[definition]{Lemma}
\newtheorem{corollary}[definition]{Corollary}

\theoremstyle{remarkbreak}
\newtheorem{remark}[definition]{Remark}
\newtheorem{example}[definition]{Example}

% Optional title block: uncomment \maketitle below to display it.
\title{PeerFS Formalization}
\author{Siddharth Mohanty}
\date{\today}


\begin{document}

\maketitle

This document formalizes the PeerFS control plane described in the README: ceilings, group
rules, topologies, mount plans and pre-flight. It does not yet formalize the data plane or the
consistency contract (Section~\ref{sec:open}).

\section{Groups and filesystems}

Let $N = \{n_1, n_2, \dots, n_k\}$ be the set of nodes in a group.

\begin{definition}[Group]
A \textbf{group} is $G = N \cup \{A\}$: the nodes $N$ together with an \textbf{authority} $A$
(the server) that holds the group rules and compiles mount plans.
\end{definition}

\begin{remark}
$A$ is a role, not necessarily a separate machine. A deployment may run the authority on one of
the nodes (for example the homelab); nothing below depends on whether it does.
\end{remark}

\begin{definition}[Paths]
Let $\Sigma$ be the set of path \emph{components}: non-empty names containing no \texttt{/},
excluding \texttt{.} and \texttt{..}. A \textbf{path} is a finite word over $\Sigma$, so
$\Path = \Sigma^*$, the empty word $\varepsilon$ is the root \texttt{/}, and $p \cdot q$ is
concatenation. Write $p \preceq q$ when $p$ is a prefix of $q$ (``$q$ lies at or under $p$'').
\end{definition}

\begin{remark}
$\preceq$ compares components, not characters. With $c = \texttt{/share}$, the path
\texttt{/shared} is \emph{not} under $c$, although the string \texttt{"/share"} is a prefix of
\texttt{"/shared"}. Implementations must check containment component-wise.
\end{remark}

\begin{lemma}\label{lem:prefix-chain}
$\preceq$ is a partial order on $\Path$, and the prefixes of any single path form a chain: if
$p \preceq x$ and $q \preceq x$, then $p \preceq q$ or $q \preceq p$.
\end{lemma}

\begin{proof}
Reflexivity, antisymmetry and transitivity are immediate for prefixes of words. If $p$ and $q$
are both prefixes of $x$, the shorter one is a prefix of the longer one, since both agree with
$x$ on their common length.
\end{proof}

\begin{definition}[Filesystem of a node]\label{def:fs}
Each node $n \in N$ has a \textbf{name tree} $F_n \subseteq \Path$: the set of paths that exist
on $n$. It is prefix-closed (if $q \in F_n$ and $p \preceq q$ then $p \in F_n$), so
$\varepsilon \in F_n$, and the graph $\Tree_n$ with vertices $F_n$ and an edge from $p$ to
$p \cdot a$ for each $a \in \Sigma$ is a tree rooted at $\varepsilon$. For $n \in N$ define:
\begin{enumerate}[(i)]
    \item $\dir_n(p) = \{\, p \cdot a \in F_n : a \in \Sigma \,\}$, the entries of directory $p$;
    \item $\sub_n(p) = \{\, q \in F_n : p \preceq q \,\}$, the subtree of $F_n$ rooted at $p$;
    \item $\canon_n : \Path \rightharpoonup F_n$, the partial function resolving symbolic links,
    undefined on paths that do not exist, dangle or loop. Symbolic links are leaves of
    $\Tree_n$; $\canon_n$ is the identity on paths none of whose prefixes is a link.
\end{enumerate}
\end{definition}

\begin{remark}
$\Tree_n$ is a tree of \emph{names}. Symbolic links do not add edges; they are handled by
$\canon_n$. Hard links give one file several names, possibly some inside and some outside a
ceiling. A name-based ceiling cannot prevent this: a hard link placed under an exported path
exports that file deliberately, just as copying it there would.
\end{remark}

\section{Ceilings}

\begin{definition}[Ceiling and exportable region]
Each node $n$ declares a \textbf{ceiling} $C(n) \subseteq \Path$, a finite set of paths. The
default is $C(n) = \emptyset$. The \textbf{exportable region} of $n$ is
\[
    \Reg(n) = \bigcup_{c \in C(n)} \sub_n(c).
\]
\end{definition}

\begin{remark}
With the default ceiling, $\Reg(n) = \emptyset$: a node exports nothing until it lists paths.
\end{remark}

\begin{lemma}\label{lem:reg-upward}
$\Reg(n)$ is closed downward in the tree: if $p \in \Reg(n)$, $x \in F_n$ and $p \preceq x$,
then $x \in \Reg(n)$.
\end{lemma}

\begin{proof}
Since $p \in \Reg(n)$ there is $c \in C(n)$ with $c \preceq p$. By transitivity $c \preceq x$,
and $x \in F_n$, so $x \in \sub_n(c) \subseteq \Reg(n)$.
\end{proof}

\section{Group rules}

\begin{definition}[Modes]
$\Mode = \{\Rd, \RW\}$, ordered by $\Rd \sqsubset \RW$.
\end{definition}

\begin{definition}[Export]
Let $\Name$ be a set of export names. An \textbf{export} is a tuple
$e = (\nu_e, o_e, \pi_e, g_e)$ where
\begin{enumerate}[(i)]
    \item $\nu_e \in \Name$ is its name;
    \item $o_e \in N$ is its \textbf{owner}, the node whose files it exposes;
    \item $\pi_e \in \Path$ is the exported path on $o_e$;
    \item $g_e : N \rightharpoonup \Mode$ is its \textbf{grant}, a finite partial function.
\end{enumerate}
\end{definition}

\begin{definition}[Group rules]
The \textbf{group rules} $R$ are the finite set of exports written by the administrator in
\texttt{DISK-RULES.yaml} (or evaluated from Pkl). Names in $R$ are \emph{intended} to be unique,
but $R$ is admin input and is not assumed well-formed; uniqueness is checked by pre-flight
(\PF{1}).
\end{definition}

\begin{definition}[Granted exports]
For a mounter $m \in N$, its \textbf{granted exports} are
$E_m = \{\, e \in R : m \in \dom g_e \,\}$.
\end{definition}

\section{Topologies}

Let $\Path_m$ denote the absolute, normalized local paths of mounter $m$ (with \texttt{\textasciitilde}
expanded); $\preceq$ is defined on them as on $\Path$.

\begin{definition}[Topology]
A \textbf{topology} for mounter $m$ is a pair $T_m = (\rho_m, \lambda_m)$ where
$\rho_m \in \Path_m$ is the \textbf{root} and $\lambda_m : \Name \rightharpoonup \Path_m$ is a
finite partial \textbf{placement}. The \textbf{default topology} is
$\tau_0 = (\texttt{/mnt/peerfs}, \emptyset)$.
\end{definition}

\begin{definition}[Targets]
Given $T_m$, the \textbf{target} of a granted export $e \in E_m$ is
\[
    \tgt_m(e) =
    \begin{cases}
        \lambda_m(\nu_e) & \text{if } \nu_e \in \dom \lambda_m, \\
        \rho_m \cdot \nu_e & \text{otherwise.}
    \end{cases}
\]
\end{definition}

\begin{definition}[Known topology]
The authority holds, for each $m$, either a topology $T_m$ received from $m$ or $\bot$ when $m$
was unreachable. Compilation uses $\hat T_m = T_m$ in the first case and $\hat T_m = \tau_0$ in
the second, and in the second case marks the plan \textbf{stale}.
\end{definition}

\section{Mount plans and pre-flight}

\begin{definition}[Mount plan]
The \textbf{mount plan} of $m$, compiled by $A$ from $R$ and $\hat T_m$, is the set of entries
\[
    P_m = \{\, (\nu_e,\ o_e{:}\pi_e,\ \tgt_m(e),\ g_e(m)) : e \in E_m \,\},
\]
read as (export, source, target, mode).
\end{definition}

\begin{definition}[Pre-flight]\label{def:preflight}
$P_m$ passes \textbf{pre-flight}, written $\mathrm{Valid}(P_m)$, if and only if all of the
following hold.
\begin{description}
    \item[\PF{1} (unique names)] For all $e \in E_m$ and $e' \in R$ with $e \neq e'$,
    $\nu_e \neq \nu_{e'}$.
    \item[\PF{2} (ceiling)] For all $e \in E_m$, $\canon_{o_e}(\pi_e)$ is defined and lies in
    $\Reg(o_e)$.
    \item[\PF{3} (placement only)] $\dom \lambda_m \subseteq \{\, \nu_e : e \in E_m \,\}$.
    \item[\PF{4} (targets)] The map $\tgt_m$ is injective on $E_m$, and its image is an
    antichain under $\preceq$: for $e \neq e'$ in $E_m$, neither
    $\tgt_m(e) \preceq \tgt_m(e')$ nor $\tgt_m(e') \preceq \tgt_m(e)$.
    \item[\PF{5} (disjoint exports)] For all $e \in E_m$ and $e' \in R$ with $e \neq e'$
    and $o_e = o_{e'}$, neither $\pi_e \preceq \pi_{e'}$ nor $\pi_{e'} \preceq \pi_e$.
    In addition, whenever the owner approves both exports under \PF{2}, neither
    $\canon_{o_e}(\pi_e) \preceq \canon_{o_e}(\pi_{e'})$ nor the reverse relation holds.
\end{description}
\end{definition}

\begin{remark}
\PF{2} is evaluated by the owner $o_e$ against its own $C(o_e)$, $F_{o_e}$ and $\canon_{o_e}$,
not by $A$: a node does not trust the central configuration with its ceiling. A refusal is
reported back to $A$, which then fails the affected plans. \PF{4} subsumes both
``colliding'' (equal) and ``nested'' targets, and covers default placements under $\rho_m$ as
well as explicit ones. \PF{5} compares both declared paths and approved canonical roots, so aliases
cannot conceal an overlap. A refused export has no approved canonical root and cannot occur in
any valid plan; its declared path still participates in \PF{5}. Only devices granted an
implicated export fail, including when they are not granted the other export. A lexical-only
offline preview cannot establish the canonical part of \PF{5} and is not admission.
\end{remark}

\section{Properties}

\subsection{Ceiling soundness}

Pre-flight checks the exported path; serving must also check every path reached through it.

\begin{definition}[Contained serving]\label{def:contained}
An owner $n$ serves \textbf{contained} if, for every request on export $e$ with $o_e = n$ for a
relative path $q \in \Path$, it serves the request only when $\canon_n(\pi_e \cdot q)$ is
defined and $\canon_n(\pi_e) \preceq \canon_n(\pi_e \cdot q)$.
\end{definition}

\begin{theorem}[Ceiling soundness]\label{thm:ceiling}
If $\mathrm{Valid}(P_m)$ and every owner serves contained, then every file $m$ reaches through
an entry of $P_m$ lies in the exportable region of that entry's owner.
\end{theorem}

\begin{proof}
Let the request be on $e \in E_m$ for relative path $q$, and let
$x = \canon_{o_e}(\pi_e \cdot q)$. By containment $x$ is defined, so $x \in F_{o_e}$, and
$\canon_{o_e}(\pi_e) \preceq x$. By \PF{2}, $\canon_{o_e}(\pi_e) \in \Reg(o_e)$. By
Lemma~\ref{lem:reg-upward}, $x \in \Reg(o_e)$.
\end{proof}

\begin{remark}
Containment is deliberately stronger than ``$x \in \Reg(o_e)$''. Consider a link inside
the \texttt{media} export that points into \texttt{/home/user/share}. Its target stays inside the
region, but following it would let a holder of \texttt{media} reach \texttt{notes} under
\texttt{media}'s grant. Containment keeps each export inside its own subtree.
\end{remark}

\subsection{Grants are exact; topologies only place}

\begin{theorem}[Grant faithfulness]\label{thm:grant}
For every topology $\hat T_m$, the projection of $P_m$ onto (export, source, mode) is
\[
    \{\, (\nu_e,\ o_e{:}\pi_e,\ g_e(m)) : e \in E_m \,\},
\]
which depends only on $R$. In particular, no topology adds an export to $P_m$ or raises a mode
above the grant.
\end{theorem}

\begin{proof}
In the definition of $P_m$, $\hat T_m$ occurs only through $\tgt_m$, which is the target
coordinate. The index set $E_m$ and the other three coordinates are functions of $R$ and $m$
alone.
\end{proof}

\begin{corollary}
Every entry of $P_m$ has mode exactly $g_e(m)$; a device granted $\Rd$ is never planned $\RW$.
\end{corollary}

\subsection{The local namespace is well-defined}

\begin{theorem}[Unambiguous resolution]\label{thm:resolve}
If $P_m$ satisfies \PF{4}, then every local path $x \in \Path_m$ lies under the target of at
most one entry: there is at most one $e \in E_m$ with $\tgt_m(e) \preceq x$.
\end{theorem}

\begin{proof}
Suppose $e \neq e'$ with $\tgt_m(e) \preceq x$ and $\tgt_m(e') \preceq x$. By
Lemma~\ref{lem:prefix-chain} the two targets are comparable. If they are equal, $\tgt_m$ is not
injective; otherwise one is a proper prefix of the other and the image is not an antichain.
Either way \PF{4} fails.
\end{proof}

Under \PF{4}, then, $m$ resolves a local path $x$ by finding the unique $e$ with
$\tgt_m(e) \preceq x$, writing $x = \tgt_m(e) \cdot q$, and requesting $q$ on $e$.

\subsection{Failure isolation}

\begin{theorem}[Isolation]\label{thm:isolation}
$\mathrm{Valid}(P_m)$ depends only on $R$, on $\hat T_m$, and on $(F_n, C(n), \canon_n)$ for
the owners $n$ of exports in $E_m$. In particular it does not depend on $\hat T_{m'}$ for any
$m' \neq m$, so a bad topology fails only its own device's plan.
\end{theorem}

\begin{proof}
By inspection of Definition~\ref{def:preflight}: \PF{1} reads $R$ and $E_m$; \PF{2} reads the
owners' filesystems and ceilings; \PF{3} and \PF{4} read $\hat T_m$ and $E_m$; \PF{5} reads
$R$ and the approval results for exports on the same owners as $E_m$. None reads another
mounter's topology.
\end{proof}

\begin{corollary}
Let $e \neq e'$ in $R$ with $\nu_e = \nu_{e'}$. The plans failing \PF{1} because of this pair
are exactly those of the devices in $\dom g_e \cup \dom g_{e'}$. Likewise, an export $e$ whose
path is outside its owner's ceiling fails \PF{2} in exactly the plans of $\dom g_e$.
\end{corollary}

\begin{proof}
$e \in E_m$ if and only if $m \in \dom g_e$, and \PF{1} and \PF{2} quantify over $E_m$.
\end{proof}

\subsection{Admission}

Let the rules and topologies change over time, and write $R^t$, $T^t_m$ and $P^t_m$ for their
values at time $t$.

\begin{definition}[Admission]
Device $m$ is \textbf{admitted} to mount at time $t$ if it sends its topology $T^t_m$ to $A$, and
$A$ compiles $P^t_m$ from $R^t$ and $T^t_m$ (not from $\tau_0$), and $\mathrm{Valid}(P^t_m)$
holds. Plans are not cached: if $A$ is unreachable, $m$ is not admitted.
\end{definition}

\begin{proposition}[Fresh admission]\label{prop:fresh}
If $m$ is admitted at time $t$, the mounts it establishes are exactly the entries of $P^t_m$,
with modes $g^t_e(m)$ from the rules current at $t$ and targets from the topology $m$ actually
holds. A stale plan is never mounted.
\end{proposition}

\begin{proof}
Admission requires a plan compiled at $t$ from $R^t$ and from $T^t_m$ itself, so the stale flag
is never set on an admitted plan. The modes follow from Theorem~\ref{thm:grant} applied to
$R^t$.
\end{proof}

\begin{remark}[Admission alone does not revoke live mounts]
Proposition~\ref{prop:fresh} constrains admission only. If $g_e(m)$ is lowered or removed at
some $t' > t$, a mount admitted at $t$ is unaffected by the definitions so far. Refusing cached
plans stops a device from \emph{re-mounting} with revoked access, not from \emph{keeping} a live
mount. One way to close the gap: number rule versions, have $A$ push the current
$g_e$ to each owner, and have the owner check $m \in \dom g_e$ and the requested operation
against $g_e(m)$ on every request, as it already checks the ceiling. Then no request after the
owner learns of the change exceeds the current grant, with exposure bounded by the propagation
delay from $A$ to the owner.
\end{remark}

\begin{remark}[Overlapping exports]
PeerFS chooses disjoint exports rather than per-path modes. \PF{5} forbids equal or nested
roots on the same owner, including approved canonical aliases. Different owners may use the
same path strings because they name different filesystems. As with \PF{2}, canonical roots
are supplied by the owner, not guessed by the authority. Pre-flight is a point-in-time check;
race-safe contained serving remains a separate data-plane obligation.
\end{remark}

\section{Implemented read-only v1 runtime}

The Linux v1 implements file RPCs over TCP carried by system Tailscale. Each endpoint resolves
the actual socket peer address through the local Tailscale WhoIs API and matches its StableID
and address against an administrator-owned enrollment registry. Device names in request bodies
are not authentication. The exporter runs as a non-root account and presents only files that
account can read; the mounter presents local UID/GID ownership. Although the control-plane
compiler preserves grants exactly, the v1 runtime intersects them with its supported operation
set: read-only. No write operation is implemented, even for an $\RW$ grant.

A successful admission creates an ephemeral session
$s = (m, v, H, d)$: authenticated mounter $m$, current persisted rule version $v$, approved
owner/root handles $H$, and expiry $d$. Sessions are not persisted. Before each file operation,
the owner authenticates the mounter and asks the authenticated authority to authorize
$(s,m,\nu_e)$. The authority permits it only if the session is unexpired, its version equals
the current rule version, the caller is the approved owner, and the current grant still
permits reading. Changed valid rule text increments the version and invalidates all sessions;
unreadable or invalid rules deny authorization. Authority or owner restart requires fresh
admission because sessions or pinned root handles have been lost.

Thus revocation is enforced at the authorization decision of each operation, rather than by
pushing grants to an owner cache. An operation already authorized before a rule change may
finish. No operation whose authorization observes the new version can use an old session.
Authority unavailability denies both admission and subsequent file operations. FUSE entry and
attribute caches are disabled and file reads use direct I/O; previously delivered application
buffers cannot be revoked. This is a conservative availability/performance tradeoff, not a
lease-based or offline consistency contract.

Ceilings and approved roots are pinned directory capabilities. Serving opens paths relative
to those handles with Linux \texttt{openat2}, using \texttt{RESOLVE\_BENEATH},
\texttt{RESOLVE\_NO\_MAGICLINKS} and \texttt{RESOLVE\_NO\_XDEV}. Absolute links, parent
traversal and nested mount crossings are refused; regular files and directories are the only
supported objects. Replacing a root pathname cannot redirect an existing handle. A directory
renamed after pinning retains that capability until expiry or revocation: the static
name-tree theorem above is not a proof of dynamic pathname-based revocation. Changing ceilings
requires restarting the exporter. File contents may change between reads; v1 makes no snapshot,
write-locking or mmap consistency promise.

The runtime is locally integration-tested with a mock identity service and actual TCP/FUSE,
including open-file revocation and symlink-swap tests. A two-machine deployment over real
Tailscale also passed binary reads, write/escape refusal, open-file revocation and automatic
restart-recovery checks. Optional user services perform explicit periodic remounts, not
transparent lease renewal; open descriptors are interrupted. These checks do not replace a
production security audit or broader load testing. See \texttt{docs/v1.md} for deployment and
protocol limits.

\section{Open questions}\label{sec:open}

\begin{enumerate}[(i)]
    \item \textbf{Consistency contract.} The README defines a mount as having ``a defined
    consistency contract'' but does not define one. The candidates (read-only exports,
    one writer per export, advisory or mandatory locks with leases) differ in what happens when
    a peer drops while holding a lock, and that is where the data-plane theorems would come
    from.
    \item \textbf{Data representation.} Copy-on-write, diffs, or whole binaries.
    \item \textbf{Access-control primitives.} Groups and \texttt{extends} would replace $g_e$'s
    domain $N$ with group expressions; Theorem~\ref{thm:grant} should survive if they are
    expanded to $N$ before compilation.
\end{enumerate}

\end{document}
